% Source: https://artofproblemsolving.com/community/c3h1707676_some_geometry_questions
% Stable identity: aops-topic-1707676-q01
Draw the line through $A$ parallel to $BC$; let it meet $EF$ at $P$ and $DC$ at $Q$. Then $ABFP$ and $ABCQ$ are parallelograms, so $PF=QC=AB=12$, hence $EP=17-12=5$ and $DQ=18-12=6$.

\begin{asy}
usepackage("amsmath");
size(7cm);
defaultpen(fontsize(10pt));
real h=sqrt(216);
pair D=(0,0),C=(18,0),A=(3,h),B=(15,h);
pair Ep=D+(2.5/15)*(A-D),F=C+(2.5/15)*(B-C);
pair Q=(6,0),P=extension(A,Q,Ep,F);
draw(A--B--C--D--cycle,linewidth(0.9));
draw(Ep--F,linewidth(0.9));
draw(A--Q,blue+dashed);
dot(Ep);dot(F);dot(P);dot(Q);
label("$A$",A,NW);label("$B$",B,NE);label("$C$",C,SE);label("$D$",D,SW);
label("$E$",Ep,W);label("$F$",F,E);label("$P$",P,NE);label("$Q$",Q,S);
label("$12$",(9,h),N);label("$18$",(12,0),S);label("$5$",(Ep+P)/2,N);label("$12$",(11,Ep.y),N);
label("\mbox{}",(min(currentpicture,user=true)+max(currentpicture,user=true))/2,black); // keeps the text after the figure black when LaTeX typesets the labels inline
\end{asy}

Since $EP\parallel DQ$, triangles $AEP$ and $ADQ$ are similar, so
\[
\dfrac{AE}{AD}=\dfrac{EP}{DQ}=\dfrac{5}{6}\ \Longrightarrow\ AE=\dfrac{5}{6}\cdot15=12.5 .
\]
Therefore $ED=15-12.5=\boxed{\dfrac{5}{2}}$.
